\section{Introduction}

The smart home is an important application domain of the Internet of Things (IoT), where automation rules connect devices, sensors, and environmental events to automate daily tasks. Most automation rules follow the Trigger-Condition-Action (TCA) model. As the number of automation rules grows, unintended rule interaction becomes more likely and may further evolve into rule conflicts, disrupting expected automation and causing security or safety risks \cite{huang2023survey,trimananda2020understanding,chi2023detecting}. For example, enabling a sprinkler upon smoke detection may indirectly trigger a leak-response rule that closes the water valve, undermining firefighting \cite{chi2023detecting}.

Recent work has further investigated rule conflict detection in IoT-based smart homes \cite{huang2023survey,huang2021conflict,li2020diac,trimananda2020understanding,yagita2015application,chaki2020fine,yu2022tapinspector,pradeep2021automating}. Existing approaches mainly include security-policy-based methods and pattern-based methods. Security-policy-based methods check whether rule execution violates predefined safety or security policies, often by modeling automation rules and checking safety properties over the resulting models \cite{celik2018soteria,celik2019iotguard,ding2021iotsafe,hsu2019safechain,sun2014conflict,ibrhim2020formal}. Pattern-based methods identify suspicious interaction structures, such as race condition, chained triggering, and conflicting actions \cite{alhanahnah2022iotcom,chi2020cross,chi2023detecting,li2020diac,xiao2019a3id,trimananda2020understanding,chaki2020fine}. Although these approaches are useful, they still have difficulty determining whether a rule interaction should be treated as a candidate rule conflict in a specific smart home setting.

This challenge arises from three factors. First, rule conflict is context-dependent. Because smart homes are heterogeneous and user preferences differ, the same rule interaction may be benign in one deployment but unsafe in another, leading to false positives or false negatives under generic detection logic. Second, many interactions propagate through physical channels with zone-dependent propagation (e.g., thermal isolation), so detectors that ignore zone boundaries may infer non-existent interactions \cite{celik2019iotguard,alhanahnah2022iotcom,chen2019multi,wang2019charting}. Third, runtime confirmation of candidate rule conflicts remains difficult. Although static analysis is comprehensive, it often over-approximates without the current system state. Existing dynamic or hybrid systems can monitor rule events and physical contexts \cite{ding2021iotsafe}, yet they do not explicitly verify whether a candidate conflict edge is actually activated at rule-trigger time.

To address these limitations, we draw three insights. Rule interaction exists not only at the logical layer but also propagates through physical channels in different zones, which requires a rule model that captures both spatial and environmental effects. Rule interaction and rule conflict should be explicitly distinguished, and the user-defined Entity Safety Configuration is used to determine whether a rule interaction should be retained as a candidate rule conflict. Finally, combining candidate conflict edges derived from static analysis with the current system state through assertion verification can reduce the over-approximation introduced by static analysis alone. Based on these observations, we propose a hybrid framework that combines offline static analysis with runtime dynamic monitoring. In the static phase, we apply formal verification to construct the Rule Interaction Dependency Graph (RIDG) and then extract the Rule Conflict Dependency Graph (RCDG) according to the user-defined Entity Safety Configuration. In the runtime phase, we use assertion verification to check whether candidate conflict edges in the RCDG are activated under the current system state, thereby identifying impending rule conflicts at runtime.

In summary, our main contributions are as follows:

\begin{itemize}
	\item We extend automation rule modeling to the Trigger-Condition-Action-Environment (TCAE) model and use the RIDG, the RCDG, and the Entity Safety Configuration to distinguish general rule interaction from candidate rule conflicts.

	\item We design a hybrid detection framework in which static analysis pre-screens interactions and limits runtime checking to candidate conflict edges instead of the full rule set.

	\item We implement a runtime mechanism based on assertion verification that combines the current system state with temporal evidence to determine whether a candidate conflict edge is activated; in our testbed, it eliminates nine false positives reported by IoTMediator while keeping runtime overhead at the millisecond level.
\end{itemize}